% LSPV-1 REPORT::SLICEcover1::2026-10-01T23:06:00-07:00
% !TEX program = xelatex
\documentclass[11pt]{article}
\usepackage[margin=0.88in]{geometry}
\usepackage{amsmath,amssymb}
\begin{document}
\section*{LSPV-1: legal-slice covering number controls the PAIR2COVER lower bound}
\textbf{STATUS: OPEN.} P versus NP remains open. No unrestricted SAT lower bound, acquisition lower bound, nonrelativizing hinge, or nonalgebrizing hinge is proved.

\paragraph{CORPUS REVIEWED.}
The current repository tree was surveyed at head. Reviewed were root \texttt{main.tex}; \texttt{AGENTS.md}, \texttt{CODEX.md}, and the pnp-research skill; all current task names with TASK-001 controlling this run; newest verified MAIN 1802 and its structural synthesis; materially relevant LSPV-1 reports through PAIRBLOCK1, including PAIR2COVER1, ORBITSLICE1, SYNCWEDGE1, DIMSUM2, RANKTRANS1 and BLOCKINNOV1; LSPV-2 through SUPPORTK1; LSPV-3 through QLEN1; LSPV-4 through WEIGHTPAD1; current governance; GROK review 2039; the Search--G and CCSPC/random-challenge branches; and referenced product-span experiment state. Repository content was surveyed for prior pair-cover, innovation, covering, sliding-window and subspace-design formulations. The automation prompt's pointer to PAIR2COVER1 as latest is stale; verified LSPV-1 reports now continue through PAIRBLOCK1.

\paragraph{LSPV-5 preflight.}
\texttt{research/main/2026-10-01\_1802\_MAIN\_1802.tex} remains the newest listed MAIN snapshot. Its verified report header, substantive OPEN body and final sentinel were inspected in this run. It keeps the synchronous exterior half OPEN and forbids promotion of structural rank to SAT runtime.

\paragraph{Primary-literature preflight.}
Shaltiel--Umans gives the algebraic extractor/orbit construction and binary conversion used by the declared SU05/Hadamard realization. van Melkebeek--Mocelin Sdroievski supplies the current RMV architecture but no theorem bounding the number of legal Hadamard slice spaces required to span the output space. Guruswami--Kopparty subspace designs control intersections of a target subspace with a family of finite-field subspaces, but their hypotheses are not established for the binary-labeled SU/vMS cyclic-window family. Thus subspace-design theory is an ANALOGY/CANDIDATE here, not a theorem transfer.

\paragraph{Source-typed object.}
For one declared source-sanctioned SU05/Hadamard realization and one legal RMV line $L$, retain ORBITSLICE1:
\[
X_s=X_{L,A^sv_0}=
\begin{bmatrix}
a_{s+1}^T\\ \vdots\\ a_{s+m}^T
\end{bmatrix},
\qquad
S_s=\operatorname{im}X_s\le\mathbb F_2^m.
\]
Let
\[
W_L=\sum_s S_s,\qquad R=\dim W_L.
\]
Define the legal slice covering number
\[
\boxed{
\tau_L=\min\left\{|Q|:Q\subseteq\mathbb Z/N\mathbb Z,\sum_{s\in Q}S_s=W_L\right\}.
}
\]
This is an intrinsic parameter of the declared legal slice family, not a runtime measure.

\paragraph{PROVED/EXACT: covering-number PAIR2COVER bound.}
For every finite legal slice family with total span dimension $R$ and covering number $\tau$, the PAIRBLOCK1 construction yields an independent direction set $T$ and legal co-occurrence edge set $E$ satisfying
\[
\boxed{
|T|=R,\qquad
|E|\ge \frac12\left(\frac{R^2}{\tau}-R\right).
}
\]
Consequently PAIR2COVER1 gives
\[
\boxed{
\rho(H)\ge
1+\frac12\left(R+\frac{R^2}{\tau}\right).
}
\]

\paragraph{Proof.}
Choose a minimum spanning subfamily $S_{s_1},\ldots,S_{s_\tau}$ and order it arbitrarily. Put
\[
W_i=\sum_{j\le i}S_{s_j},\qquad
d_i=\dim W_i-\dim W_{i-1}.
\]
Minimality implies $d_i\ge1$ after any ordering of this chosen minimal subfamily, and
\[
\sum_{i=1}^{\tau}d_i=R.
\]
PAIRBLOCK1 constructs an independent $T$ of size $R$ with at least
\[
\sum_i\binom{d_i}{2}
=\frac12\left(\sum_i d_i^2-R\right)
\]
legal edges. Cauchy--Schwarz gives
\[
\sum_i d_i^2\ge\frac{(\sum_i d_i)^2}{\tau}=\frac{R^2}{\tau}.
\]
Substitution proves both displayed bounds.

\paragraph{Immediate structural corollaries.}
If $R=\Theta(m)$ and
\[
\tau=o(m),
\]
then
\[
\rho(H)=\omega(m),
\]
so a global $O(m)$ product-span theorem is false for that SAME\_REALIZATION family. In particular, if every slice has rank at most $k$ and the total span $R=m$ can be reached using
\[
\tau=O(m/k)
\]
legal slices with $k\to\infty$, then
\[
\rho(H)=\Omega(mk).
\]
Conversely, if $R=\Theta(m)$ and an architecture-specific theorem establishes $\rho(H)=O(m)$, the covering-number inequality forces
\[
\boxed{\tau=\Omega(m).}
\]
Thus a linear product-span theorem would require the legal window spaces to be highly redundant in a precise covering sense: spanning a linear-dimensional output space must require linearly many legal slices.

\paragraph{Single-slice obstruction and why multiple slices matter.}
A single legal slice of rank $r$ already gives
\[
\rho(H)\ge1+r+\binom r2.
\]
Hence one slice would suffice only if $r=\omega(\sqrt m)$. In the declared Hadamard realization $r\le k=\lceil\log_2p\rceil$; this observation does not by itself meet the source-parameter target and cannot replace the multi-slice incidence analysis. The covering-number theorem identifies the missing collective quantity without assuming individual high rank.

\paragraph{What this does not prove.}
No upper bound on $\tau_L$ is proved for an actual vMS canonical LDE line. Full linear span $R=m$ alone is insufficient: $\tau_L$ may be $\Theta(m)$, corresponding to essentially one-dimensional innovation per indispensable slice. Generic large field, low odd-prime field degree, list decoding, locality, short source, canonicity, or subspace-design existence does not bound $\tau_L$ for this binary-labeled cyclic family.

\paragraph{PQ/cross-domain classification.}
\textbf{EXACT OBJECT/THEOREM TRANSFER:} fixed-code Schur/exterior objects, quotient-basis PAIRBLOCK1, Cauchy--Schwarz, finite-dimensional span and covering number, and the declared SU05/Hadamard orbit matrices.
\textbf{ANALOGY/CANDIDATE:} subspace designs, dimension expanders, matroid covering/packing, shortening filtrations, quadratic-relation spectra, Hilbert functions and syzygies. A source-specific theorem relating these to $\tau_L$ is missing.
\textbf{FAILED TRANSFER:} generic subspace-design existence, random-code heuristics, cryptographic hardness/random oracles, VLSI locality, control/observability, distributed cuts, topology, proof-system lower bounds, large field alone, or low field degree as automatic bounds on $\tau_L$ or unrestricted SAT time.

\paragraph{DUPLICATION / CONTRADICTION AUDIT.}
PAIR2COVER1 supplied $\rho(H)\ge1+|T|+|E|$. BLOCKINNOV1 supplied $\dim U\ge\sum_i\binom{d_i}{2}$. PAIRBLOCK1 constructed $T,E$ from an arbitrary innovation sequence. The new content is the minimization over spanning legal subfamilies and the Cauchy--Schwarz elimination of the individual $d_i$, yielding a bound solely in terms of $(R,\tau)$. No contradiction was found with ORBITSLICE1, SYNCWEDGE1, DIMSUM2, SUPPORTK1, RANKTRANS1 or the causal/search branches.

\paragraph{WHAT CHANGED RELATIVE TO ROOT MAIN.}
Root MAIN already records the squared-innovation target. This run compresses that target further to two measurable legal-family parameters:
\[
(R,\tau_L).
\]
For $R=\Theta(m)$, proving merely $\tau_L=o(m)$ suffices for superlinear product span. The next experiment need not optimize all orderings of legal slices before testing the route.

\paragraph{ADOPTED / REJECTED / RETIRED PRIOR COMPONENTS.}
Adopted: PAIR2COVER1, ORBITSLICE1, SYNCWEDGE1, DIMSUM2, BLOCKINNOV1, PAIRBLOCK1 and SAME\_REALIZATION discipline.
Rejected as necessary targets: spanning $m$-clique, maximal exterior rank, one-state high rank, and exact optimization of the innovation ordering.
Retired as generic shortcuts: short-source, canonicity, odd-prime field degree, list decoding, locality, repeated-root structure, or large field alone.

\paragraph{STALE OR OVERSTATED CLAIMS FOUND.}
The run prompt's latest-report pointer is stale. Root MAIN also does not yet contain PAIRBLOCK1 because its previous reconciliation write was blocked. The mathematical content of current root MAIN remains conservative; this run should reconcile both PAIRBLOCK1 and the new covering-number theorem into the LSPV-1 owned region only.

\paragraph{Hostile controls and barriers.}
SAME\_REALIZATION fixes $L,f,A,\kappa,m$, legal shifts and downstream consumer. Semantic/algebraic oracles can reveal every $S_s$, $R$ and $\tau_L$, so the theorem is algebraic/relativizing and not an acquisition lower bound. DEAD/COPY/DELAY preserve the extensional slice family; INTERP/RECOMP require retyping. Padding and Generic Avoid remain controls. Affine/Tseitin with Gaussian elimination, Horn, 2-SAT, bounded treewidth and independent products remain easy controls. Boolean source degree remains OPEN. Natural-proofs and algebrization barriers are not crossed.

\paragraph{Strict status and next target.}
\textbf{PROVED/EXACT:} for every declared finite legal slice family,
\[
\rho(H)\ge1+\tfrac12(R+R^2/\tau).
\]
\textbf{CONDITIONAL:} if one SAME\_REALIZATION family has $R=\Theta(m)$ and $\tau=o(m)$, then its product-span rank is $\omega(m)$.
\textbf{OPEN:} $R$ and $\tau_L$ for the earliest reconstruction-compatible actual vMS LDE line; Boolean source degree; exact author-unique bit-level invocation; charged acquisition; arbitrary-algorithm transfer; nonrelativizing/nonalgebrizing hinge; final diagonal.

The single next target is to compute or symbolically bound the legal slice covering number. For the exact cyclic windows $S_s=\operatorname{im}X_s$, greedily add a slice maximizing $\dim(S_s+W)-\dim W$ until $W=W_L$, record the achieved cover size and compare it with the true minimum on feasible small cases. Search first for legal LDE families with $R=\Theta(m)$ but $\tau=\Theta(m)$; such families are the direct hostile control.

\begin{thebibliography}{4}
\bibitem{SU05} R. Shaltiel and C. Umans, \emph{Simple Extractors for All Min-Entropies and a New Pseudorandom Generator}, J. ACM 52(2):172--216, 2005.
\bibitem{vMS25} D. van Melkebeek and N. Mocelin Sdroievski, \emph{Instance-Wise Hardness and Refutation versus Derandomization for Arthur--Merlin Protocols}, Computational Complexity 34:15, 2025.
\bibitem{GK13} V. Guruswami and S. Kopparty, \emph{Explicit Subspace Designs}, FOCS 2013, pp. 608--617; ECCC TR13-060.
\bibitem{Rand} H. Randriambololona, \emph{On products and powers of linear codes under componentwise multiplication}, Contemporary Mathematics 637:3--78, 2015.
\end{thebibliography}
\end{document}
% END_OF_VERIFIED_REPORT::LSPV-1::SLICECOVER1
